\documentclass[10pt]{article}
\usepackage[T1]{fontenc}
\usepackage[utf8]{inputenc}
\usepackage[french]{babel}
\usepackage{amsmath,amssymb,mathrsfs,mathpazo}
\usepackage{geometry}
\geometry{a4paper,margin=2cm}
\usepackage{graphicx,array,tabularx,tabularray,colortbl,multirow,multicol,booktabs,enumitem,pifont,ulem,fancybox,diagbox,cancel}
\usepackage{pstricks,pstricks-add,pst-plot,pst-tree,pst-node,pst-text,pst-func,pst-math,pst-bspline,pst-3dplot,pst-eucl,pst-solides3d}
\let\pscalculate\undefined
\usepackage{pst-bezier}
\usepackage{tikz,pgfplots,tkz-tab}
\usetikzlibrary{arrows,arrows.meta,patterns,calc,positioning,trees,decorations.markings,shapes.geometric,shapes.misc,shapes.arrows,intersections,decorations.pathmorphing,decorations.pathreplacing,backgrounds,shadows,babel}
\pgfplotsset{compat=1.18}
\usepgfplotslibrary{fillbetween}
\usepackage[np]{numprint}
\usepackage{listings,makecell,pas-tableur,twoopt}
\let\barre\undefined
\usepackage[pstricks]{tablvar}
\definecolor{darkgreen}{RGB}{0,100,0}
\providecommand{\red}{\color{red}}
\providecommand{\blue}{\color{blue}}
\providecommand{\green}{\color{green}}
\providecommand{\black}{\color{black}}
\providecommand{\white}{\color{white}}
\providecommand{\gray}{\color{gray}}
\providecommand{\cyan}{\color{cyan}}
\providecommand{\magenta}{\color{magenta}}
\providecommand{\yellow}{\color{yellow}}
\providecommand{\violet}{\color{violet}}
\providecommand{\orange}{\color{orange}}
\providecommand{\euro}{\texteuro}
\providecommand{\EUR}{\texteuro}
\providecommand{\ieme}{\textsuperscript{e}}
\providecommand{\degres}{\ensuremath{^\circ}}
\providecommand{\degree}{\ensuremath{^\circ}}
\providecommand{\vv}[1]{\overrightarrow{#1}}
\providecommand{\Coord}[1]{#1}
\providecommand{\decosix}{}
\providecommand{\starredbullet}{\ensuremath{\bullet}}
\pagestyle{empty}
\setlength{\parindent}{0pt}
\usepackage{adjustbox}\DeclareUnicodeCharacter{2236}{:}\DeclareUnicodeCharacter{2190}{\ensuremath{\leftarrow}}
\providecommand{\R}{\mathbb{R}}
\providecommand{\N}{\mathbb{N}}
\providecommand{\Z}{\mathbb{Z}}
\providecommand{\Q}{\mathbb{Q}}
\providecommand{\D}{\mathbb{D}}
\providecommand{\C}{\mathbb{C}}
\providecommand{\e}{{\mathrm{e}}}
\providecommand{\vect}[1]{\overrightarrow{#1}}
\providecommand{\vectt}[1]{\overrightarrow{\mathrm{#1}}}
\providecommand{\vvt}[1]{\overrightarrow{\mathrm{#1}}}
\providecommand{\np}[2][]{{#2~#1}}
\providecommand{\vv}[1]{{\overrightarrow{#1}}}
\renewcommand{\d}{{\mathrm{d}}}
\providecommand{\degres}{{^\circ}}
\providecommand{\degre}{{^\circ}}
\providecommand{\blue}{\color{blue}}
\providecommand{\red}{\color{red}}
\providecommand{\enskip}{\quad}
\providecommand{\llbracket}{\lbrack\!\lbrack}
\providecommand{\rrbracket}{\rbrack\!\rbrack}
\providecommand{\euro}{{\text{€}}}
\providecommand{\no}{{\text{n°}}}
\providecommand{\interval}[2]{{[#1;#2]}}
\providecommand{\intervFO}[2]{{[#1;#2[}}
\providecommand{\dotfill}{\ldots}

\providecommand{\ds}{\displaystyle}
\providecommand{\textding}[1]{#1}
\providecommand{\Oij}{\(\mathrm{O};\vec{\imath},\vec{\jmath}\)}
\providecommand{\Oijk}{\(\mathrm{O};\vec{\imath},\vec{\jmath},\vec{k}\)}
\providecommand{\Ouv}{\(\mathrm{O};\vec{u},\vec{v}\)}
\providecommand{\intervFF}[2]{\ensuremath{[#1;#2]}}
\renewcommand{\euro}{\ifmmode\text{\texteuro}\else\texteuro\fi}
\DeclareUnicodeCharacter{00B0}{\ifmmode^{\circ}\else\textdegree\fi}
\DeclareUnicodeCharacter{20AC}{\ifmmode\text{\texteuro}\else\texteuro\fi}
% Énoncés : Éducation nationale. Transcription et corrigés : APMEP, Denis Vergès et contributeurs. Découpage et classement : CoopMaths, Sébastien Lozano, Éric Elter et contributeurs. Adaptation de la présentation : Maths Explorer.
% CC BY-SA : https://coopmaths.fr/www/about/
% https://coopmaths.fr/alea/static/bac/2021/tex/bac_2021_09_metropole_7_cor.tex
\begin{document}
\noindent\textit{APMEP / CoopMaths. CC BY-SA. Présentation adaptée par Maths Explorer.}\par\medskip
\textbf{Commun à tous les candidats}


















\begin{enumerate}
\item 








Pour $t = - 2$, on trouve les coordonnées de B.
\item 
Le vecteur $\vect{\text{AB}}$ admet pour coordonnées : $\begin{pmatrix}2 - 1~;~1 - 0~;~0 - 2\end{pmatrix}$ soit $\begin{pmatrix}1~;~1~;~- 2\end{pmatrix}$








\item 









$M(x~;~y~;~z) \in (\text{AB})$ s'il existe $t \in \R$ tel que $\vect{\text{A}M} = t\vect{\text{AB}}$ soit :

$\left\{\begin{array}{l !{=} l}
x - 1& 1\times t\\
y - 0& 1\times t\\
z - 2& - 2\times t
\end{array}\right.\, t \in \R \iff 
\left\{\begin{array}{l !{=} r}
x & 1 +  t\\
y &t\\
z & 2 - 2t
\end{array}\right.\, t \in \R 
$.

En posant $t = 1 - u$, on obtient :

$M(x~;~y~;~z) \in (\text{AB}) \iff \left\{\begin{array}{l !{=} r}
x & 2 - u\\
y &1 - u\\
z & 2u
\end{array}\right.\, u \in \R$. Donc réponse B.
\item 







La droite $\Delta$ a pour vecteur directeur $\delta\begin{pmatrix}2\\1\\-1\end{pmatrix}$

Soit $\mathcal{P}$ le plan dont on cherche une équation.

$M(x~;~y~;~z) \in \mathcal{P} \iff \vect{\text{C}M} \cdot \delta = 0 \iff 2(x - 0) + 1(y - 1) -1(z - 2) = 0 \iff$

$ 2x + y - 1 - z + 2 = 0 \iff 2x + y - z + 1 = 0$.
\item 







En faisant apparaître le point A dans chaque vecteur (Chasles), on obtient :

$\vect{\text{OA}} + \vect{\text{AD}}  = 3\vect{\text{OA}}  - \vect{\text{OA}}  - \vect{\text{AB}}  - \vect{\text{OA}}  - \vect{\text{AC}} \iff \vect{\text{AD}}  = - \vect{\text{AB}} - \vect{\text{AC}} $ : le vecteur $\vect{\text{AD}}$ est une combinaison linéaire des vecteurs $\vect{\text{AB}}$ et $\vect{\text{AC}}$, ces trois vecteurs sont donc coplanaires.
\end{enumerate}

\bigskip

\end{document}
